\subsection{MSVC: x86}

\lstinputlisting{patterns/04_scanf/2_global/ex2_MSVC.asm}

\RU{В целом ничего особенного. Теперь \TT{x} объявлена в сегменте \TT{\_DATA}. 
Память для неё в стеке более не выделяется.
Все обращения к ней происходит не через стек, а уже напрямую. 
Неинициализированные глобальные переменные не занимают места в исполняемом файле
(и действительно, зачем в исполняемом файле
нужно выделять место под изначально нулевые переменные?), но тогда, когда к этому месту в памяти
кто-то обратится, \ac{OS} подставит туда блок, состоящий из нулей\footnote{Так работает \ac{VM}}.}
\EN{In this case the \TT{x} variable is defined in the \TT{\_DATA} segment and no memory is allocated in the local stack. It is accessed directly, not through the stack. 
Uninitialized global variables take no space in the executable file
(indeed, why one needs to allocate space for variables initially set to zero?), 
but when someone accesses their address, 
the \ac{OS} will allocate a block of zeroes there\footnote{That is how a \ac{VM} behaves}.}
\ES{En este caso, la variable \TT{x} está definida en el segmento \TT{\_DATA} y no se asigna memoria en la pila local. Se accede a ella directamente, no a través de la pila. Las variables globales no inicializadas no ocupan espacio en el archivo ejecutable (de hecho, ¿por qué se necesitaría asignar espacio para variables establecidas inicialmente a cero?), pero cuando alguien accede a su dirección, el \ac{OS} asignará un bloque de ceros allí\footnote{Así es como se comporta una \ac{VM}}.}

\RU{Попробуем изменить объявление этой переменной:}
\EN{Now let's explicitly assign a value to the variable:}
\ES{Ahora asignemos explícitamente un valor a la variable:}

\lstinputlisting{patterns/04_scanf/2_global/default_value.c.\LANG}

\RU{Выйдет в итоге:}\EN{We got:}\ES{Obtenemos:}

\begin{lstlisting}
_DATA	SEGMENT
_x	DD	0aH

...
\end{lstlisting}

\RU{Здесь уже по месту этой переменной записано \TT{0xA} с типом DD (dword = 32 бита).}
\EN{Here we see a value \TT{0xA} of DWORD type (DD stands for DWORD = 32 bit) for this variable.}
\ES{Aquí vemos un valor \TT{0xA} de tipo DWORD (DD significa DWORD = 32 bits) para esta variable.}

\RU{Если вы откроете скомпилированный .exe-файл в \IDA, то увидите, что \IT{x} 
находится в начале сегмента \TT{\_DATA}, после этой переменной будут текстовые строки.}
\EN{If you open the compiled .exe in \IDA, you can see the \IT{x} variable placed at the beginning of 
the \TT{\_DATA} segment, and after it you can see text strings.}
\ES{Si abre el ejecutable compilado en \IDA, podrá ver la variable \IT{x} ubicada al principio del segmento \TT{\_DATA}, y después de ella podrá ver las cadenas de texto.}

\RU{А вот если вы откроете в \IDA .exe скомпилированный в прошлом примере, 
где значение \IT{x} не определено, то вы увидите:}
\EN{If you open the compiled .exe from the previous example in \IDA, where the value of \IT{x} was not
set, you would see something like this:}
\ES{Si abre en \IDA el ejecutable compilado del ejemplo anterior, donde no se estableció el valor de \IT{x}, verá algo como esto:}

\begin{lstlisting}
.data:0040FA80 _x              dd ?                    ; DATA XREF: _main+10
.data:0040FA80                                         ; _main+22
.data:0040FA84 dword_40FA84    dd ?                    ; DATA XREF: _memset+1E
.data:0040FA84                                         ; unknown_libname_1+28
.data:0040FA88 dword_40FA88    dd ?                    ; DATA XREF: ___sbh_find_block+5
.data:0040FA88                                         ; ___sbh_free_block+2BC
.data:0040FA8C ; LPVOID lpMem
.data:0040FA8C lpMem           dd ?                    ; DATA XREF: ___sbh_find_block+B
.data:0040FA8C                                         ; ___sbh_free_block+2CA
.data:0040FA90 dword_40FA90    dd ?                    ; DATA XREF: _V6_HeapAlloc+13
.data:0040FA90                                         ; __calloc_impl+72
.data:0040FA94 dword_40FA94    dd ?                    ; DATA XREF: ___sbh_free_block+2FE
\end{lstlisting}

\RU{\TT{\_x} обозначен как \TT{?}, наряду с другими переменными не требующими инициализации. 
Это означает, что при загрузке .exe в память, место под всё это выделено будет и будет заполнено
нулевыми байтами \cite[6.7.8p10]{C99TC3}. 
Но в самом .exe ничего этого нет. Неинициализированные переменные не занимают места в исполняемых файлах. 
Это удобно для больших массивов, например.}
\EN{\TT{\_x} is marked with \TT{?} with the rest of the variables that do not need to be initialized. 
This implies that after loading the .exe to the memory, a space for all these variables is to be 
allocated and filled with zeroes \cite[6.7.8p10]{C99TC3}. 
But in the .exe file these uninitialized variables do not occupy anything.
This is convenient for large arrays, for example.}
\ES{\TT{\_x} está marcada con \TT{?}, junto con el resto de las variables que no necesitan ser inicializadas. Esto implica que después de cargar el archivo ejecutable en la memoria, se asignará un espacio para todas estas variables y se llenará con ceros \cite[6.7.8p10]{C99TC3}. Pero en el archivo ejecutable estas variables no inicializadas no ocupan nada. Esto resulta muy conveniente para matrices grandes, por ejemplo.}

\ifdefined\IncludeOlly
\clearpage
\subsection{MSVC: x86 + \olly}
\index{\olly}

\RU{Тут даже проще}\EN{Things are even simpler here}\ES{Las cosas son aún más sencillas aquí}:

\begin{figure}[H]
\centering
\includegraphics[scale=\FigScale]{patterns/04_scanf/2_global/ex2_olly_1.png}
\caption{\olly: \RU{после исполнения \scanf}\EN{after \scanf execution}\ES{después de la ejecución de \scanf}}
\label{fig:scanf_ex2_olly_1}
\end{figure}

\RU{Переменная хранится в сегменте данных.}
\EN{The variable is located in the data segment.}
\ES{La variable se encuentra en el segmento de datos.}
\RU{Кстати, после исполнения инструкции \PUSH (заталкивающей адрес $x$) адрес появится в стеке, 
и на этом элементе можно нажать правой кнопкой, выбрать \q{Follow in dump}.}
\EN{After the \PUSH instruction (pushing the address of $x$) gets executed, 
the address appears in the stack window. Right-click on that row and select \q{Follow in dump}.}
\ES{Por cierto, después de que se ejecute la instrucción \PUSH (que introduce la dirección de $x$), la dirección aparece en la ventana de la pila. Haga clic derecho en esa fila y seleccione \q{Follow in dump}.}
\RU{И в окне памяти слева появится эта переменная.}
\EN{The variable will appear in the memory window on the left.}
\ES{La variable aparecerá en la ventana de memoria a la izquierda.}

\RU{После того как в консоли введем 123, здесь появится}\EN{After we have entered 123 in the console,} 
\TT{0x7B}\EN{ appears in the memory window (see the highlighted screenshot regions)}\ES{ aparece en la ventana de memoria (vea las regiones resaltadas de la captura de pantalla)}.

\RU{Почему самый первый байт это}\EN{But why is the first byte} \TT{7B}?
\RU{По логике вещей, здесь должно было бы быть}\EN{Thinking logically,} \TT{00 00 00 7B}\EN{ should be
there}\ES{ debería estar allí}.
\RU{Это называется}\EN{The cause for this is referred as } \gls{endianness}, \RU{и в x86 принят формат }\EN{and x86 uses }\IT{little-endian}.
\ES{La causa de esto se conoce como} \gls{endianness}, \ES{y x86 utiliza} \IT{little-endian}.
\RU{Это означает, что в начале записывается самый младший байт, а заканчивается самым старшим байтом}%
\EN{This implies that the lowest byte is written first, and the highest written last}.
\ES{Esto implica que el byte más bajo se escribe primero y el más alto al final.}
\RU{Больше об этом}\EN{Read more about it at}\ES{Lea más sobre esto en}: \myref{sec:endianness}.

\RU{Позже из этого места в памяти 32-битное значение загружается в \EAX и передается в}
\EN{Back to the example, the 32-bit value is loaded from this memory address into \EAX and passed to} \printf.
\ES{Volviendo al ejemplo, el valor de 32 bits se carga desde esta dirección de memoria en \EAX y se pasa a} \printf.

\RU{Адрес переменной $x$ в памяти}\EN{The memory address of $x$ is}\ES{La dirección de memoria de $x$ es} \TT{0x00C53394}.

\clearpage
\RU{В \olly{} мы можем посмотреть карту памяти процесса (Alt-M) и увидим, что этот адрес
внутри PE-сегмента \TT{.data} нашей программы}%
\EN{In \olly we can review the process memory map (Alt-M)
and we can see that this address is inside the \TT{.data} PE-segment of our program}:
\ES{En \olly podemos revisar el mapa de memoria del proceso (Alt-M) y podemos ver que esta dirección está dentro del segmento PE \TT{.data} de nuestro programa}:

\begin{figure}[H]
\centering
\includegraphics[scale=\FigScale]{patterns/04_scanf/2_global/ex2_olly_2.png}
\caption{\olly: \RU{карта памяти процесса}\EN{process memory map}\ES{mapa de memoria del proceso}}
\label{fig:scanf_ex2_olly_2}
\end{figure}

\fi

\ifdefined\IncludeGCC
\subsection{GCC: x86}

\index{ELF}
\RU{В Linux всё почти также. За исключением того, что если значение \TT{x} не определено, 
то эта переменная будет находится в сегменте \TT{\_bss}.
В \ac{ELF} этот сегмент имеет такие атрибуты:}
\EN{The picture in Linux is near the same, with the difference that the uninitialized variables are located in the \TT{\_bss} segment. 
In \ac{ELF} file this segment has the following attributes:}
\ES{El panorama en Linux es casi el mismo, con la diferencia de que las variables no inicializadas se ubican en el segmento \TT{\_bss}. En el archivo \ac{ELF}, este segmento tiene los siguientes atributos:}

\begin{lstlisting}
; Segment type: Uninitialized
; Segment permissions: Read/Write
\end{lstlisting}

\RU{Ну а если сделать статическое присвоение этой переменной какого-либо
значения, например, 10, то она будет находится 
в сегменте \TT{\_data},
это сегмент с такими атрибутами:}
\EN{If you, however, initialise the variable with some value e.g. 10, 
it is to be placed in the \TT{\_data} segment, which has the following attributes:}
\ES{Sin embargo, si inicializa la variable con algún valor, por ejemplo, 10, se colocará en el segmento \TT{\_data}, que tiene los siguientes atributos:}

\begin{lstlisting}
; Segment type: Pure data
; Segment permissions: Read/Write
\end{lstlisting}
\fi

\subsection{MSVC: x64}

\lstinputlisting[caption=MSVC 2012 x64]{patterns/04_scanf/2_global/ex2_MSVC_x64.asm.\LANG}

\RU{Почти такой же код как и в}\EN{The code is almost the same as in}\ES{El código es casi el mismo que en} x86.
\RU{Обратите внимание что для \TT{scanf()} адрес переменной $x$ передается
при помощи инструкции \LEA, а во второй \printf передается само значение переменной при помощи \MOV}
\EN{Please note that the address of the $x$ variable is passed to \TT{scanf()} using a \LEA instruction,
while the variable's value is passed to the second \printf using a \MOV instruction}.
\ES{Tenga en cuenta que la dirección de la variable $x$ se pasa a \TT{scanf()} mediante una instrucción \LEA, mientras que el valor de la variable se pasa al segundo \printf mediante una instrucción \MOV.}
\TT{DWORD PTR}\RU{\EMDASH{}это часть языка ассемблера (не имеющая отношения к машинным кодам) 
показывающая, что тип переменной в памяти именно 32-битный, 
и инструкция \MOV должна быть здесь закодирована соответственно}\EN{ is a part of the assembly language (no relation to the machine code), indicating that the variable data size is 32-bit and the \MOV instruction
has to be encoded accordingly}.
\ES{ es una parte del lenguaje ensamblador (sin relación con el código de máquina), que indica que el tamaño de los datos de la variable es de 32 bits y la instrucción \MOV debe codificarse en consecuencia.}
